%% RESUMO NA LINGUA VERNACULA (elemento obrigatorio -- NBR 6028)
%% Paragrafo unico, 150 a 500 palavras: objetivo, metodo, resultados, conclusao.
%% Termina com as palavras-chave (ja vem de \palavraschave em config/dados.tex).

\begin{resumo}
A inteligência artificial já está mudando a trajetória dos trabalhadores
brasileiros? Esta dissertação investiga se a exposição ocupacional à
inteligência artificial está associada a mudanças nos padrões de transição no
mercado de trabalho após o lançamento público do ChatGPT, em 30 de novembro de
2022. Utiliza-se o painel rotativo da PNAD Contínua entre 2019 e 2025, com uma
amostra analítica de 3.963.359 pessoas-transições com exposição atribuída à
ocupação de origem, distribuídas em 413 códigos ocupacionais e 27 trimestres. A
exposição é mensurada por um índice de exposição ocupacional à inteligência
artificial, originalmente construído para a classificação ocupacional
norte-americana e compatibilizado com a brasileira. A estratégia empírica adota
um desenho de diferenças em diferenças com tratamento contínuo, estimado por
modelo de probabilidade linear ponderado, com efeitos fixos de ocupação de
origem e de unidade da federação por trimestre e com a teletrabalhabilidade
interagida com o período. Após novembro de 2022, observa-se aumento generalizado
das taxas medidas de mobilidade ocupacional e, além desse movimento, um
gradiente associado à exposição à inteligência artificial: cada unidade
adicional do índice está associada a uma variação entre os dois períodos 2,19
pontos percentuais maior na transição para ocupações menos expostas e 1,36 ponto
percentual menor na transição para ocupações mais expostas. O resultado mais
surpreendente aparece nas transições relacionadas ao vínculo de trabalho: maior
exposição associa-se a menor saída do emprego e a menor passagem do emprego
formal para o informal, contrariando uma interpretação simples de substituição
tecnológica. Em conjunto, os resultados são compatíveis com um ajuste que
aparece na recomposição das trajetórias ocupacionais, e não na eliminação de
postos de trabalho, ainda que o desenho não isole o canal responsável por esse
padrão. As estimativas devem ser interpretadas como associações condicionais, e
não como efeitos causais: o trabalho reporta um único cenário de estimação, sem
bateria de sensibilidade, e a série utilizada atravessa uma descontinuidade na
taxa medida de mudança de ocupação, coincidente com a alteração no modo de
coleta da pesquisa em 2020 e 2021. A principal contribuição do estudo é oferecer
evidência brasileira, em larga escala, para um debate ainda fortemente apoiado
em estudos internacionais, com implicações para as políticas públicas de
trabalho, qualificação profissional e proteção social.
\end{resumo}
